\documentclass{beamer}

\input{MathMacros}

\usepackage{beamerthemesplit}
\usepackage{enumerate}
\usepackage{amssymb}
\usepackage[all,cmtip]{xy}
\usepackage[mathscr]{eucal}
\usepackage[natbib=true,style=authoryear,backend=bibtex,useprefix=true]{biblatex}
\addbibresource{reading.bib}
\usepackage{graphicx,color}

\title{Foundations of Large Language Models: What They Are and How They Behave}
\author{Reuben Brasher}
\date{\today}

\begin{document}

\frame{\titlepage}

\section[Outline]{}
\frame{\tableofcontents}

\section{Lecture 1}

\frame
{
   \frametitle{A Familiar Shift in a New Domain}

    We are experiencing a familiar technological shift in a new cognitive domain.

    \begin{itemize}
        \item Calculators reduced the need for manual arithmetic but not mathematical understanding.
        \item Students now ask the same question in a new form: why learn this if AI can do it.
        \item The domain of external assistance has shifted from computation to language and reasoning.
        \item External tools can perform tasks without replacing the need for internal competence.
        \item Real-world situations still require unassisted thinking, speaking, and judgment.
        \item The analogy is useful only if we understand what LLMs actually are and are not.
    \end{itemize}
}

\frame
{
   \frametitle{What LLMs Fundamentally Are}

    Large language models are best understood as systems for sequence transformation.

    \begin{itemize}
        \item At a fundamental level, LLMs map sequences of tokens to new sequences.
        \item They are trained on large corpora such as Common Crawl and other aggregated text sources.
        \item Their objective is next-token prediction rather than truth or correctness.
        \item Outputs are conditioned entirely on the provided context and prior generated tokens.
        \item They do not contain structured knowledge in the manner of databases or symbolic systems.
        \item Their behavior is better described as transformation than as understanding or reasoning.
    \end{itemize}
}

\frame
{
   \frametitle{How Context Shapes Output}

    The structure and content of the input largely determine the direction of the output.

    \begin{itemize}
        \item The model responds to patterns in the prompt rather than to an internal world model.
        \item Small changes in phrasing or framing can substantially alter the generated response.
        \item Structured input constrains the model toward more stable and predictable outputs.
        \item Ungrounded prompts invite unconstrained sampling from the training distribution.
        \item Context functions simultaneously as instruction and as constraint on generation.
        \item Effective use requires deliberate construction of input rather than passive querying.
    \end{itemize}
}

\frame
{
   \frametitle{Why LLMs Perform Well}

    LLMs are effective when tasks can be framed as transformations between representations.

    \begin{itemize}
        \item Strong performance appears when translating between formats such as text, schema, and code.
        \item The models excel at recognizing and reproducing common structural patterns.
        \item Tasks with clear mappings between representations are especially reliable.
        \item Performance depends on statistical regularities learned from large training corpora.
        \item Structured prompts reduce ambiguity and improve consistency of the output.
        \item This explains why some tasks appear intelligent while others degrade very quickly.
    \end{itemize}
}

\frame
{
   \frametitle{Evidence from Structured Tasks}

    Empirical evidence shows that reliability increases when work is staged through explicit structure.

    \begin{itemize}
        \item In structured cases such as status effects, schema to JSON to code transitions were accurate.
        \item Consistent field definitions enabled alignment across multiple representations.
        \item Unit tests generated from structured inputs were valid and executable.
        \item Minimal correction was needed when scaffolding and constraints were present.
        \item The model maintained coherence across several transformation steps.
        \item Reliability therefore appears to be conditional on structure rather than inherent capability.
    \end{itemize}
}

\frame
{
   \frametitle{Failure Mode 1: Hallucination and Fabrication}

    When grounding is weak, LLMs often produce plausible statements that are false.

    \begin{itemize}
        \item Models may invent facts, citations, or numerical values that appear credible.
        \item In the EDA experiment, trends were inferred from minimal data and presented as factual.
        \item Fabricated regression test outputs were produced without any actual computation.
        \item Hallucinations are patterned and plausible rather than purely random.
        \item Users often accept such outputs because fluency is mistaken for reliability.
        \item This creates systematic risk in research, analysis, and decision-making work.
    \end{itemize}
}

\frame
{
   \frametitle{Failure Mode 2: Structural vs Semantic Errors}

    LLMs often preserve formal structure while losing semantic correctness.

    \begin{itemize}
        \item Code may compile or appear correct while implementing the wrong logic.
        \item In the dice example, probabilistic reasoning was replaced by a uniform distribution.
        \item Structural patterns are easier for the model to reproduce than domain-specific meaning.
        \item The model can preserve form without validating function.
        \item Such errors often remain hidden unless explicit tests are applied.
        \item Correct structure therefore does not imply correct semantics.
    \end{itemize}
}

\frame
{
   \frametitle{Failure Mode 3: Context Drift}

    Long interactions often degrade because context accumulates noise and irrelevant patterning.

    \begin{itemize}
        \item Extended conversations accumulate unintended context that affects later outputs.
        \item Models may introduce irrelevant artifacts such as code blocks in unrelated tasks.
        \item Output style and format may drift unless actively corrected.
        \item Earlier errors can propagate and distort future responses.
        \item The model does not maintain clean topic separation unless the user enforces it.
        \item Managing context is therefore an active responsibility rather than a passive assumption.
    \end{itemize}
}

\frame
{
   \frametitle{The Core Limitation}

    The central limitation of LLMs is that they generate by constrained pattern completion rather than truth-tracking.

    \begin{itemize}
        \item There is no internal mechanism that independently verifies correctness.
        \item The model cannot on its own validate facts, proofs, or computations.
        \item It lacks persistent state and stable internal representations across interactions.
        \item It does not reliably track contradictions unless prompted and constrained to do so.
        \item Reliability must therefore be imposed externally through structure and validation.
        \item This limitation is fundamental to the system rather than a temporary defect.
    \end{itemize}
}

\frame
{
   \frametitle{Implications for Users}

    Effective use of LLMs requires disciplined control over inputs, outputs, and validation.

    \begin{itemize}
        \item Begin with grounded inputs such as documents, data, or explicit references.
        \item Treat outputs as drafts or hypotheses rather than as final answers.
        \item Correct the model explicitly when it deviates from the task.
        \item Avoid copying generated text without careful inspection and revision.
        \item Use iteration, comparison, and cross-checking to improve reliability.
        \item The model interprets silence as acceptance and reinforces that interaction pattern.
    \end{itemize}
}

\frame
{
   \frametitle{Discussion}

    Understanding LLM behavior leads directly to the problem of disciplined use.

    \begin{itemize}
        \item LLMs are powerful systems but they are inherently unreliable when unconstrained.
        \item Their strongest performance depends on structure, staging, and explicit constraints.
        \item Unstructured use leads predictably to degraded outputs and poor reasoning outcomes.
        \item Effective use is therefore a matter of interaction design rather than mere access.
        \item The real challenge is not possession of the tool but disciplined integration of it.
        \item This leads to the next question: how can we use these systems without losing rigor.
    \end{itemize}
}

\begin{frame}[t,allowframebreaks]
\frametitle{References}
\printbibliography
\end{frame}

\end{document}
